\documentclass[11pt]{article}
\usepackage[margin=1in]{geometry}
\usepackage{graphicx,amsmath,amssymb,booktabs,hyperref,xcolor}
\graphicspath{{/tmp/cheby_h0/figs/}}
\hypersetup{colorlinks=true,linkcolor=blue,urlcolor=blue}
\title{Does the tabulated-$\mu$ architecture really need Chebyshev?\\
       a three-way head-to-head}
\author{Fixed-Point Factory --- experiment note}
\date{June 2026}

\begin{document}
\maketitle

\begin{abstract}
\noindent
The tabulated-$\mu$ refactoring (\texttt{tab\_explain.pdf}) reduces the
per-cell contour cost to a single 1D table lookup, with Phase-B clearing
running independently at every cube cell. Once that simplification is in
place, the natural question is: does the underlying conjecture $P$ still
need to be a tensor-Chebyshev polynomial, or could a simpler CDF-uniform
linear-grid representation do the same job? We answer empirically by
running three operator variants at $N{=}7$, $\tau{=}\gamma{=}1$:
\begin{itemize}
  \item[(A)] \textbf{Cheb tab} --- Lobatto--atanh nodes, atanh-stretched
        GL quadrature in $\xi_a$, Chebyshev polynomial slice evaluation.
  \item[(B)] \textbf{Linear-on-Lobatto} --- \emph{same nodes, same
        quadrature} as A, but slice evaluation by 2D linear
        interpolation. Isolates the effect of the interpolant alone.
  \item[(C)] \textbf{Linear uniform-$u$} --- uniform-in-$u$ nodes on
        $[-U_{\max}, U_{\max}]$, uniform-$u$ GL quadrature, linear
        interpolation. Drops the Chebyshev infrastructure entirely.
\end{itemize}
Headline: \textbf{A's floor is 10$\times$ better than B's, and 35$\times$
better than C's.} Variant C's floor does \emph{not} improve as you refine
$G_{\rm lin}$ --- the uniform-$u$ GL quadrature is the wrong measure for
$\int f_v\,du$ and the iteration converges to a different (wrong) FP.
Variant B is the most informative: same nodes and same quadrature as A,
so the entire 10$\times$ gap is attributable to losing the smooth Chebyshev
derivative $\partial P/\partial u_b$ in the co-area integrand. The
verdict: \textbf{at the current resolution, Chebyshev is buying you one
order of magnitude in residual quality} --- expensive to abandon at
$N{=}7$, harder to justify if the operator's structural floor were much
lower.
\end{abstract}

\section{Design of the experiment}

The three operators all share the tabulated-$\mu$ architecture from
\texttt{tab\_explain.pdf}: build $\mu(p, u_k)$ on a $G_p\!\times\!G$
table once per $\Phi$, then look up at every cube cell and run CRRA
clearing. They differ only in how the underlying $P$ is represented:

\begin{center}
\begin{tabular}{lcccc}
\toprule
Operator & nodes for $P$ & GL quad in $\xi_a$? & slice eval & root in $\xi_b$ \\
\midrule
A.\ Cheb tab           & Lobatto+atanh & atanh-stretched & Cheb polynomial & bisection on poly \\
B.\ Lin-on-Lobatto     & Lobatto+atanh & atanh-stretched & 2D bilinear     & sign-change on grid \\
C.\ Linear uniform-$u$ & uniform $u$   & uniform $u$     & 2D bilinear     & sign-change on grid \\
\bottomrule
\end{tabular}
\end{center}

Both linear variants store $P$ as raw values at the grid nodes; no
coefficient transform. The grid spacings are $G_p{=}51$ for the
$p$-table in all three; the cube-clearing grid is $G{=}7$ for A and B,
and we sweep $G_{\rm lin}\in\{9,13,17,21,25,33\}$ for C.

Starting conjecture: $P_0(u) = \sigma(0.5\,\tau\sum_k u_k)$ on each
operator's own grid. Iteration: Anderson type-II with memory $m{=}8$,
80 iterations.

\section{Result 1: the floor under iteration}

\begin{center}
\includegraphics[width=0.97\linewidth]{three_ops_01_floor.png}
\end{center}

\begin{tabular}{lcc}
\toprule
operator & min $\|F\|_\infty$ (80 Anderson iters) & ratio vs A \\
\midrule
A. Cheb tab            & $3.04\!\times\!10^{-3}$ & $1\times$ \\
B. Linear on Lobatto   & $2.97\!\times\!10^{-2}$ & $9.8\times$ worse \\
C. Linear uniform-$u$  & $1.12\!\times\!10^{-1}$ & $37\times$ worse \\
\bottomrule
\end{tabular}

\vspace{0.5em}
A reaches the known structural floor at $N{=}7$ ($3\!\times\!10^{-3}$,
diagnosed in \texttt{fp\_analysis.pdf}). B and C both undershoot it
substantially. The key comparison is \textbf{A vs B}: identical nodes,
identical quadrature, only the interpolation differs. The full 10$\times$
gap is attributable to losing the smooth $\partial P/\partial u_b$ in the
co-area integrand.

\paragraph{Why piecewise-linear interpolation hurts.} The co-area
integrand carries a $1/|\partial P/\partial u_b|$ factor. With Chebyshev
$P$, this derivative is a smooth polynomial (\texttt{chebder} coefficient
recurrence is exact). With piecewise-linear $P$, $\partial P/\partial u_b$
is \emph{piecewise constant} --- it jumps at every Lobatto grid line.
Those jumps inject high-frequency noise into the $\mu$ table, which feeds
back through clearing and limits how close the iteration can drive the
residual.

\section{Result 2: variant C does not converge with refinement}

The natural defense of C (linear uniform-$u$) is that refining $G_{\rm
lin}$ should converge to the truth. It does not:

\begin{center}
\includegraphics[width=0.97\linewidth]{three_ops_04_lin_sweep.png}
\end{center}

\textbf{Left.} As $G_{\rm lin}$ grows from 9 to 33, the floor first
\emph{worsens} (from $2.6\!\times\!10^{-2}$ at $G{=}9$ to
$1.6\!\times\!10^{-1}$ at $G{=}33$). It never reaches the Cheb floor.
\textbf{Right.} Wall time grows like $G_{\rm lin}^3$ (driven by the
per-cube-cell CRRA bisection), so refinement is cubic-expensive.

\textbf{Why refinement fails:} The Hellwig co-area integral involves
$\int f_v(u_a)\,f_v(u_b(u_a))\,/\,|\partial P/\partial u_b|\,du_a$
where $f_v$ is a Gaussian density centred at $u = \pm \tfrac{1}{2}$ with
unit standard deviation (at $\tau{=}1$). The Cheb operator uses
atanh-stretched GL nodes, which concentrate quadrature points near
$u\approx 0$ where $f_v\cdot f_v$ is largest. Variant C uses GL nodes
uniformly mapped to $[-U_{\max},U_{\max}]$, putting most nodes in the
\emph{tails} where $f_v$ is small --- starving the integral of the
mass it needs. Refining the storage grid does not fix this: the GL
points are still in the wrong places. The result is a systematically
biased $\mu$ table and a different (wrong) fixed point, with slope
$\alpha^* \approx 0.66$ instead of A's $0.34$ at $\gamma{=}1$.

\section{Result 3: wall time}

\begin{center}
\includegraphics[width=0.85\linewidth]{three_ops_02_walltime.png}
\end{center}

\begin{tabular}{lcc}
\toprule
operator & wall time per $\Phi$ & floor \\
\midrule
A. Cheb tab            & $9.5$ ms & $3.0\!\times\!10^{-3}$ \\
B. Linear on Lobatto   & $5.2$ ms ($\mathbf{1.8\times}$ faster) & $3.0\!\times\!10^{-2}$ \\
C. Linear uniform-$u$ ($G_{\rm lin}{=}21$) & $123$ ms (13$\times$ slower) & $1.1\!\times\!10^{-1}$ \\
\bottomrule
\end{tabular}

\vspace{0.5em}
Variant B \emph{is} faster than A --- by $\sim 2\times$ --- because the
slice evaluation is a simple bilinear blend instead of a Chebyshev value
recurrence, and the root extraction is a single sign-change scan rather
than bisection on the polynomial. But the floor is 10$\times$ worse.

\textbf{Cost-of-quality:} B gives you 2$\times$ faster $\Phi$ for 10$\times$
worse residual. If the goal is just to get an FP at $\|F\|_\infty\sim 3
\!\times\!10^{-2}$ (perhaps for a coarse phase diagram), B is the right
choice. If the goal is the structural floor at $\sim 10^{-3}$, you need
A.

Variant C's 13$\times$ slowdown comes from $G_{\rm lin}{=}21$ giving
$21^3 = 9261$ cube cells (vs A's $7^3 = 343$); the per-cell CRRA bisection
dominates and there's no recovery from refinement.

\section{Result 4: $\gamma$-dependence}

\begin{center}
\includegraphics[width=0.85\linewidth]{three_ops_03_gamma.png}
\end{center}

The three operators agree on the qualitative shape (slope rising with
$\gamma$, as expected from CRRA risk-aversion concentrating demand on
the more probable state), but the \emph{level} is systematically off:

\begin{tabular}{cccc}
\toprule
$\gamma$ & A. Cheb & B. Lin-Lob & C. Lin-uniform \\
\midrule
0.3 & 0.305 & 0.485 & 0.539 \\
1.0 & 0.342 & 0.596 & 0.658 \\
5.0 & 0.463 & 0.731 & 0.667 \\
\bottomrule
\end{tabular}

\vspace{0.5em}
B and C both over-estimate the slope by 50--90\% across the range. The
discrepancy comes from the same source as the residual floor: the
mis-estimated $\partial P/\partial u_b$ in the co-area integrand and (for
C) the mis-placed quadrature nodes corrupt the $\mu(p, u_k)$ table.

\section{Result 5: slice visualisation}

A direct look at $P$ after one $\Phi$ application (slice fixing $u_2 = $
middle Lobatto node):

\begin{center}
\includegraphics[width=\linewidth]{three_ops_05_slices.png}
\end{center}

A and B look qualitatively similar (the underlying conjecture is
$\sigma(0.5T)$, smooth and antisymmetric about $T{=}0$), with A's slice
slightly sharper at the boundaries due to the smooth Cheb derivative.
C looks markedly different: the slope across the diagonal is shallower
and the boundary regions are more diffuse, consistent with its
under-resolved tail quadrature.

\section{Verdict: what role is Chebyshev still playing?}

\begin{enumerate}
  \item \textbf{Storing $P$ as a tensor-Chebyshev polynomial} ---
        \emph{not essential.} Phase B (clearing) only reads $P$ at the
        cube grid points; any representation that gives the same node
        values works.
  \item \textbf{Lobatto--atanh node placement + atanh-stretched GL
        quadrature} --- \emph{essential.} Even with linear interpolation
        of $P$ (variant B), keeping the Lobatto--atanh structure gives
        $3\!\times\!10^{-2}$. Replacing it with uniform-$u$ (variant C)
        loses another 4$\times$ in floor and never recovers with
        refinement.
  \item \textbf{Smooth $\partial P/\partial u_b$ via \texttt{chebder}} ---
        \emph{worth one order of magnitude in floor.} The whole gap from
        $3\!\times\!10^{-3}$ (A) to $3\!\times\!10^{-2}$ (B) is paid for
        by the smooth derivative in the co-area integrand. This is the
        single largest contributor.
\end{enumerate}

\paragraph{Practical implication.} The structural floor at $N{=}7$ is
$\sim 1.7\!\times\!10^{-3}$ (LM+Newton, \texttt{fp\_analysis.pdf}).
Variant A reaches $3.0\!\times\!10^{-3}$ --- close to that floor.
Variant B reaches $3.0\!\times\!10^{-2}$ --- a factor of 17$\times$
above the structural floor. So Chebyshev is currently the bottleneck
between "achieves the architecture's floor" and "leaves 17$\times$ of
room on the table". \emph{If we were already near a much higher floor
(say $10^{-1}$), linear-on-Lobatto would be the cheaper choice.}

\paragraph{What would close the gap?} A C$^1$ piecewise interpolant
(PCHIP / monotone cubic Hermite) on Lobatto nodes would give a smooth
$\partial P/\partial u_b$ without going to full Chebyshev. Predicted
floor: somewhere between B ($3\!\times\!10^{-2}$) and A
($3\!\times\!10^{-3}$). This is the natural next experiment if speed
matters more than the last factor of three.

\paragraph{Summary recommendation.} Keep Chebyshev for now. The smooth
derivative is genuinely doing work, not just costing time. The only
component to consider replacing is the storage layer (Cheb coeffs $\to$
raw values) once $N$ grows large enough that the coefficient
representation becomes a memory burden --- which is not the regime
we're in at $N \lesssim 12$.

\section*{Reproducibility}

Figures generated by \texttt{plot\_three\_ops.py}; raw numbers in
\texttt{three\_ops.json}. Operator sources:
\texttt{cheby\_numba\_tab.py} (A), \texttt{cheby\_linterp\_on\_lob.py}
(B), \texttt{cheby\_linear\_op.py} (C). All on branch
\texttt{claude/study-fixed-point-economics-y12PB}.

\end{document}
